\documentclass[a4paper,12pt]{article}

% ==============================================================================
% PAQUETES DE CONFIGURACIÓN Y TIPOGRAFÍA
% ==============================================================================
\usepackage[utf8]{inputenc}
\usepackage[T1]{fontenc}
\usepackage[spanish,es-nodecimaldot,es-noindentfirst]{babel}
\usepackage{amsmath,amssymb}
\usepackage{graphicx}
\usepackage{float}
\usepackage{xcolor}
\usepackage{colortbl}
\usepackage{geometry}
\usepackage{fancyhdr}
\usepackage{booktabs}
\usepackage{tabularx}
\usepackage{array}
\newcolumntype{Y}{>{\raggedright\arraybackslash}X}
\newcolumntype{C}{>{\centering\arraybackslash}X}
\usepackage{enumitem}
\usepackage{listings}
\usepackage{microtype}
\usepackage{titlesec}
\usepackage{tikz}
\usetikzlibrary{shapes.geometric, arrows.meta, positioning, calc, fit, backgrounds, shadows}
\usepackage[most]{tcolorbox}
\usepackage{hyperref}

% ==============================================================================
% DIMENSIONES DE PÁGINA (GEOMETRY)
% ==============================================================================
\geometry{
    top=2.5cm,
    bottom=2.5cm,
    left=2.5cm,
    right=2.5cm,
    headheight=25pt,
    footskip=1.5cm
}

% ==============================================================================
% PALETA INSTITUCIONAL (DESIGN SYSTEM WAYKA / GACETA SUCRE)
% ==============================================================================
\definecolor{colorSucre}{HTML}{800000}       % Rojo Sucre / Garnet Primario
\definecolor{colorSucreDark}{HTML}{550000}   % Rojo Sucre Oscuro
\definecolor{colorAzul}{HTML}{1B365D}        % Azul Institucional / Navy Secundario
\definecolor{colorAzulDark}{HTML}{132744}    % Azul Institucional Oscuro
\definecolor{colorAzulLight}{HTML}{2A4F85}   % Azul Claro Acento
\definecolor{colorGrisFondo}{HTML}{F4F6F9}   % Gris de Fondo / Superficie
\definecolor{colorBorde}{HTML}{E2E8F0}       % Gris de Bordes / Separadores
\definecolor{colorTexto}{HTML}{212529}       % Gris Oscuro para Texto Principal
\definecolor{colorTextoMuted}{HTML}{6C757D}  % Gris Neutro Secundario
\definecolor{colorExito}{HTML}{1E7E34}       % Verde Éxito
\definecolor{colorExitoBg}{HTML}{E6F4EA}     % Verde Fondo
\definecolor{colorAlerta}{HTML}{B71C1C}      % Rojo Alerta
\definecolor{colorAlertaBg}{HTML}{FCE8E6}    % Rojo Fondo
\definecolor{colorWarning}{HTML}{8D6500}     % Amarillo Advertencia
\definecolor{colorWarningBg}{HTML}{FFF8E1}   % Amarillo Fondo
\definecolor{colorInfo}{HTML}{006064}        % Cian Informativo
\definecolor{colorInfoBg}{HTML}{E0F7FA}      % Cian Fondo

% ==============================================================================
% CONFIGURACIÓN DE HIPERVÍNCULOS (HYPERREF)
% ==============================================================================
\hypersetup{
    colorlinks=true,
    linkcolor=colorSucre,
    citecolor=colorAzul,
    urlcolor=colorAzul,
    pdftitle={Informe Técnico Sprint 1 - Sistema Wayka MVP},
    pdfauthor={Pablo - Tech Lead \& Senior Software Architect},
    pdfsubject={Gestión Documental y Workflow Institucional},
    pdfkeywords={Wayka, Sprint 1, JWT, RBAC, MySQL, REST API, Sucre, GAMS}
}

% ==============================================================================
% CONFIGURACIÓN DE CÓDIGO FUENTE (LISTINGS)
% ==============================================================================
\lstdefinelanguage{json}{
    basicstyle=\ttfamily\footnotesize\color{colorTexto},
    numbers=left,
    numberstyle=\tiny\color{colorTextoMuted},
    stepnumber=1,
    numbersep=8pt,
    showstringspaces=false,
    breaklines=true,
    breakatwhitespace=true,
    frame=single,
    rulecolor=\color{colorBorde},
    backgroundcolor=\color{colorGrisFondo},
    stringstyle=\color{colorSucre}\bfseries,
    keywordstyle=\color{colorAzul}\bfseries,
    commentstyle=\color{colorTextoMuted}\itshape,
    morestring=[b]",
    literate=
      {á}{{\'{a}}}{1}
      {é}{{\'{e}}}{1}
      {í}{{\'{i}}}{1}
      {ó}{{\'{o}}}{1}
      {ú}{{\'{u}}}{1}
      {ñ}{{\~{n}}}{1}
      {:}{{{\color{colorAzulDark}{:}}}}{1}
      {,}{{{\color{colorAzulDark}{,}}}}{1}
      {\{}{{{\color{colorSucreDark}{\{}}}}{1}
      {\}}{{{\color{colorSucreDark}{\}}}}}{1}
      {[}{{{\color{colorSucreDark}{[}}}}{1}
      {]}{{{\color{colorSucreDark}{]}}}}{1},
}

\lstdefinestyle{sqlstyle}{
    language=SQL,
    basicstyle=\ttfamily\footnotesize\color{colorTexto},
    numbers=left,
    numberstyle=\tiny\color{colorTextoMuted},
    stepnumber=1,
    numbersep=8pt,
    showstringspaces=false,
    breaklines=true,
    frame=single,
    rulecolor=\color{colorBorde},
    backgroundcolor=\color{colorGrisFondo},
    keywordstyle=\color{colorAzul}\bfseries,
    stringstyle=\color{colorSucre},
    commentstyle=\color{colorTextoMuted}\itshape
}

% ==============================================================================
% ENCABEZADOS Y PIES DE PÁGINA (FANCYHDR)
% ==============================================================================
\pagestyle{fancy}
\fancyhf{}
\renewcommand{\headrulewidth}{0.8pt}
\renewcommand{\footrulewidth}{0.5pt}

\renewcommand{\headrule}{\hbox to\headwidth{\color{colorSucre}\leaders\hrule height \headrulewidth\hfill}}
\renewcommand{\footrule}{\hbox to\headwidth{\color{colorBorde}\leaders\hrule height \footrulewidth\hfill}}

\fancyhead[L]{\small\bfseries\color{colorAzul}SISTEMA WAYKA MVP \textbar{} GESTIÓN DOCUMENTAL}
\fancyfoot[L]{\footnotesize\color{colorTextoMuted}Gobierno Autónomo Municipal de Sucre -- Práctica Laboral SHC170}
\fancyfoot[R]{\small\bfseries\color{colorAzul}Página \thepage}

% ==============================================================================
% FORMATO DE TÍTULOS Y SECCIONES (TITLESEC)
% ==============================================================================
\titleformat{\section}
  {\normalfont\Large\bfseries\color{colorSucreDark}}
  {\thesection.}{0.8em}{}[\color{colorSucre}\titlerule]

\titleformat{\subsection}
  {\normalfont\large\bfseries\raggedright\color{colorAzul}}
  {\thesubsection.}{0.6em}{}

\titleformat{\subsubsection}
  {\normalfont\normalsize\bfseries\color{colorTexto}}
  {\thesubsubsection.}{0.5em}{}

\setlength{\parskip}{0.6em}
\setlength{\parindent}{0pt}
\setlength{\emergencystretch}{2em}

% ==============================================================================
% DOCUMENTO PRINCIPAL
% ==============================================================================
\begin{document}
\hypersetup{pageanchor=false}

% ------------------------------------------------------------------------------
% CARÁTULA / PORTADA INSTITUCIONAL
% ------------------------------------------------------------------------------
\begin{titlepage}
    \begin{tikzpicture}[remember picture, overlay]
        % Franja lateral decorativa izquierda institucional
        \fill[colorSucre] (current page.north west) rectangle ($(current page.north west) + (1.2cm, -\paperheight)$);
        \fill[colorAzul] ($(current page.north west) + (1.2cm, 0)$) rectangle ($(current page.north west) + (1.5cm, -\paperheight)$);
        
        % Barra decorativa superior de encabezado
        \fill[colorGrisFondo] ($(current page.north west) + (1.5cm, 0)$) rectangle ($(current page.north east) + (0, -3.5cm)$);
        \draw[colorBorde, line width=1.5pt] ($(current page.north west) + (1.5cm, -3.5cm)$) -- (current page.north east);
    \end{tikzpicture}

    \vspace{-1.5cm}
    \begin{flushright}
        {\large\bfseries\color{colorAzul}GOBIERNO AUTÓNOMO MUNICIPAL DE SUCRE}\\[0.15cm]
        {\footnotesize\color{colorSucre}\textbf{PROYECTO WAYKA: SISTEMA DE GESTIÓN DOCUMENTAL Y WORKFLOW}}
    \end{flushright}

    \vspace{3.5cm}

    \begin{center}
        \begingroup
            \color{colorSucreDark}
            \Huge\textbf{INFORME TÉCNICO DE INGENIERÍA}\\[0.3cm]
            \LARGE\textbf{SPRINT 1: ARQUITECTURA BASE Y FUNDACIÓN}
        \endgroup

        \vspace{0.8cm}
        \noindent\makebox[\linewidth]{\color{colorAzul}\rule{0.75\linewidth}{2pt}}
        \vspace{0.8cm}

        {\large\textbf{\color{colorTexto}Consolidación del Modelo Relacional, Motor de Migraciones DDL,}\\[0.15cm]
        \textbf{\color{colorTexto}Autenticación Criptográfica JWT y Matriz de Autorización RBAC}}

        \vspace{3cm}

        \begin{tcolorbox}[
            colback=colorGrisFondo,
            colframe=colorAzul,
            arc=3mm,
            width=0.9\linewidth,
            boxrule=1.2pt
        ]
            \centering
            \vspace{0.2cm}
            \begin{tabularx}{\linewidth}{@{}lY@{}}
                \textbf{\color{colorAzul}Producto:} & Sistema Wayka MVP (Versión 1.0.0-sprint1) \\
                \textbf{\color{colorAzul}Fase / Sprint:} & Sprint 1 (Semanas 1 y 2 -- Fundación Técnica) \\
                \textbf{\color{colorAzul}Período de Ejecución:} & 01 de Septiembre al 14 de Septiembre de 2026 \\
                \textbf{\color{colorAzul}Asignatura:} & SHC170 -- Práctica Laboral \\
                \textbf{\color{colorAzul}Estado de Ejecución:} & \textbf{\color{colorExito}Completada}
            \end{tabularx}
            \vspace{0.2cm}
        \end{tcolorbox}
    \end{center}

    \vfill

    \begin{center}
        \small\color{colorTextoMuted}
        Sucre, Bolivia \textbar{} Gestión 2026
    \end{center}
\end{titlepage}
\hypersetup{pageanchor=true}

% ------------------------------------------------------------------------------
% TABLA DE CONTENIDOS
% ------------------------------------------------------------------------------
\newpage
\tableofcontents
\newpage

% ==============================================================================
% SECCIÓN 1: RESUMEN EJECUTIVO
% ==============================================================================
\section{Resumen Ejecutivo}

El presente documento constituye el Informe Técnico Exhaustivo correspondiente a la culminación del \textbf{Sprint 1 (Fundación)} del proyecto \textbf{Sistema Wayka MVP (Sistema de Gestión Documental y Workflow Institucional)}, implementado para modernizar la administración y trazabilidad de expedientes y correspondencias en el Gobierno Autónomo Municipal de Sucre.

Durante este primer ciclo de dos semanas (01 al 14 de Septiembre de 2026), se establecieron con éxito los cimientos arquitectónicos, estructurales y de seguridad de la plataforma, garantizando una base desacoplada, escalable y con alta cohesión técnica.

\subsection{Hitos y Logros Principales}
\begin{itemize}[leftmargin=1.5em]
    \item \textbf{Capa de Persistencia y Motor DDL:} Diseño e implementación del esquema relacional unificado en MySQL 8.0+ mediante un motor automatizado de migraciones versionadas (\texttt{migrator.js}) y 11 scripts incrementales ordenados cronológicamente (\texttt{001\_create\_personas.sql} a \texttt{011\_create\_parametros.sql}).
    \item \textbf{Subsistema Criptográfico y Autenticación JWT:} Construcción de los servicios de emisión, validación y renovación de tokens criptográficos \emph{JSON Web Token} (JWT) con firma HMAC-SHA256, incorporación de secretos distribuidos y hashing de contraseñas de alta seguridad mediante \texttt{bcrypt} (10 rondas de salt).
    \item \textbf{Autorización Multi-Rol y Modelo RBAC:} Estructuración de la matriz de control de acceso basada en roles (\emph{Role-Based Access Control} -- RBAC) para los cuatro roles esenciales del MVP: \emph{Administrador de Sistema}, \emph{Administrador de Wayka}, \emph{Ventanilla Única} y \emph{Funcionario}, soportando pertenencia simultánea a múltiples oficinas orgánicas.
    \item \textbf{Capa de Controladores y Servicios REST API:} Despliegue de los endpoints de autenticación, control de perfiles, conmutación dinámica de contexto operativo (\texttt{switch-role}) y modelos de acceso a datos (DAO) para Personas, Usuarios, Ubicaciones Orgánicas y Asignaciones.
    \item \textbf{Garantía de Calidad y Pruebas Automatizadas:} Verificación exhaustiva con suites de pruebas unitarias y de integración end-to-end (\texttt{testAuth.js} y \texttt{testDataModel.js}), logrando una tasa de éxito del 100\% (17/17 aserciones críticas aprobadas sin regresiones).
\end{itemize}

% ==============================================================================
% SECCIÓN 2: ARQUITECTURA DE SOFTWARE Y MODELO DE DATOS
% ==============================================================================
\section{Arquitectura de Software y Modelo de Datos}

El Sistema Wayka está diseñado bajo un enfoque de \textbf{Arquitectura en Capas Limpia (Clean Layered Architecture)} y desacoplamiento API-First. La capa de persistencia se apoya en un motor relacional MySQL 8.0+ con motor de almacenamiento InnoDB, garantizando cumplimiento riguroso de las propiedades ACID (Atomicidad, Consistencia, Aislamiento y Durabilidad).

\subsection{Estrategia de Modelado y Persistencia Relacional}
La arquitectura de datos del Sprint 1 resuelve tres requerimientos críticos de la gestión pública:
\begin{enumerate}[leftmargin=1.5em]
    \item \textbf{Unificación de la Identidad Civil e Institucional:} Separación explícita entre la identidad física de un ciudadano o servidor público (\texttt{personas}) y su cuenta digital de acceso (\texttt{usuarios}), permitiendo que una persona mantenga su historial aun cuando cambie de cargo o dependencia.
    \item \textbf{Organigrama Jerárquico Recursivo:} La tabla \texttt{ubicaciones\_org} implementa una estructura de árbol adyacente (\emph{Adjacency List}) mediante el campo autorreferencial \texttt{padre\_id}, calculando automáticamente los niveles administrativos ($N+1$).
    \item \textbf{Modelo Multi-Rol Contextualizado:} A través de la tabla asociativa \texttt{usuario\_roles}, un usuario puede ejercer diferentes funciones según la oficina u órgano municipal en el que esté asignado, gobernado por restricciones de unicidad compuesta (\texttt{uk\_\allowbreak usuario\_\allowbreak rol\_\allowbreak ubicacion}) y fechas de caducidad.
\end{enumerate}

\subsection{Estrategia de Migraciones y Auditoría de Datos}
Para evitar discrepancias entre entornos de desarrollo, pruebas y producción, se prescindió del modelado manual ad-hoc, adoptando un subsistema de migraciones idempotente gobernado por la tabla de control \texttt{\_migrations}:
\begin{itemize}[leftmargin=1.5em]
    \item \textbf{Integridad Referencial Estricta:} Todas las llaves foráneas se definieron con la cláusula \texttt{ON DELETE RESTRICT}, bloqueando eliminaciones en cascada accidentales que pudieran comprometer la trazabilidad de expedientes oficiales.
    \item \textbf{Borrado Lógico (\emph{Soft Delete}):} Las entidades transaccionales y de catálogo implementan el flag booleano \texttt{activo DEFAULT 1}. La capa de lógica de negocio valida dependencias antes de desactivar cualquier registro (ej. una persona no puede desactivarse si cuenta con un usuario activo).
    \item {\raggedright\textbf{Auditoría Temporal Nativa:} Inclusión sistemática de marcas temporales \texttt{created\_at TIMESTAMP DEFAULT CURRENT\_\allowbreak TIMESTAMP} y \texttt{updated\_at TIMESTAMP DEFAULT CURRENT\_\allowbreak TIMESTAMP ON UPDATE CURRENT\_\allowbreak TIMESTAMP}.\par}
\end{itemize}

\subsection{Diagrama Entidad-Relación (ERD) en TikZ}
La Figura~\ref{fig:erd_sprint1} ilustra la topología relacional de las entidades fundacionales construidas durante el Sprint 1.

\begin{figure}[H]
\centering
\resizebox{\linewidth}{!}{%
\begin{tikzpicture}[
    scale=0.82, transform shape,
    tablebox/.style={
        rectangle,
        draw=black,
        text=black,
        line width=0.6pt,
        fill=white,
        inner sep=0pt
    },
    rel/.style={
        -{Stealth[scale=1.0]},
        line width=0.6pt,
        color=black
    },
    card/.style={
        font=\tiny\bfseries\color{black},
        fill=white,
        inner sep=1pt
    }
]

    % NODO: PERSONAS
    \node (personas) at (0, 7.5) [tablebox] {
        \begin{tabular}{>{\raggedright\arraybackslash}p{5.0cm}}
            \centering\arraybackslash\textbf{\footnotesize personas} \\
            \hline
            \scriptsize\ttfamily \textbf{PK} id : INT (AI)\\
            \scriptsize\ttfamily \textbf{UQ} ci : VARCHAR(25)\\
            \scriptsize\ttfamily \phantom{XX} nombres : VARCHAR(100)\\
            \scriptsize\ttfamily \phantom{XX} apellido\_paterno : VARCHAR(100)\\
            \scriptsize\ttfamily \phantom{XX} apellido\_materno : VARCHAR(100)\\
            \scriptsize\ttfamily \phantom{XX} email : VARCHAR(120)\\
            \scriptsize\ttfamily \phantom{XX} telefono : VARCHAR(30)\\
            \scriptsize\ttfamily \phantom{XX} activo : BOOLEAN
        \end{tabular}
    };

    % NODO: USUARIOS
    \node (usuarios) at (0, 1.2) [tablebox] {
        \begin{tabular}{>{\raggedright\arraybackslash}p{5.0cm}}
            \centering\arraybackslash\textbf{\footnotesize usuarios} \\
            \hline
            \scriptsize\ttfamily \textbf{PK} id : INT (AI)\\
            \scriptsize\ttfamily \textbf{FK} persona\_id : INT\\
            \scriptsize\ttfamily \textbf{UQ} login : VARCHAR(50)\\
            \scriptsize\ttfamily \phantom{XX} password\_hash : VARCHAR(255)\\
            \scriptsize\ttfamily \phantom{XX} cargo : VARCHAR(120)\\
            \scriptsize\ttfamily \phantom{XX} ultimo\_acceso : DATETIME\\
            \scriptsize\ttfamily \phantom{XX} activo : BOOLEAN
        \end{tabular}
    };

    % NODO: USUARIO_ROLES (Intermedia)
    \node (userroles) at (7.5, 1.2) [tablebox] {
        \begin{tabular}{>{\raggedright\arraybackslash}p{5.2cm}}
            \centering\arraybackslash\textbf{\footnotesize usuario\_roles} \\
            \hline
            \scriptsize\ttfamily \textbf{PK} id : INT (AI)\\
            \scriptsize\ttfamily \textbf{FK} usuario\_id : INT\\
            \scriptsize\ttfamily \textbf{FK} rol\_id : INT\\
            \scriptsize\ttfamily \textbf{FK} ubicacion\_org\_id : INT\\
            \scriptsize\ttfamily \phantom{XX} nivel\_acceso : VARCHAR(50)\\
            \scriptsize\ttfamily \phantom{XX} fecha\_expiracion : DATE\\
            \scriptsize\ttfamily \phantom{XX} es\_principal : BOOLEAN\\
            \scriptsize\ttfamily \phantom{XX} activo : BOOLEAN
        \end{tabular}
    };

    % NODO: ROLES
    \node (roles) at (7.5, 7.5) [tablebox] {
        \begin{tabular}{>{\raggedright\arraybackslash}p{5.2cm}}
            \centering\arraybackslash\textbf{\footnotesize roles} \\
            \hline
            \scriptsize\ttfamily \textbf{PK} id : INT (AI)\\
            \scriptsize\ttfamily \textbf{UQ} codigo : VARCHAR(50)\\
            \scriptsize\ttfamily \phantom{XX} nombre : VARCHAR(100)\\
            \scriptsize\ttfamily \phantom{XX} descripcion : TEXT\\
            \scriptsize\ttfamily \phantom{XX} activo : BOOLEAN
        \end{tabular}
    };

    % NODO: UBICACIONES_ORG
    \node (ubicaciones) at (14.5, 1.2) [tablebox] {
        \begin{tabular}{>{\raggedright\arraybackslash}p{5.2cm}}
            \centering\arraybackslash\textbf{\footnotesize ubicaciones\_org} \\
            \hline
            \scriptsize\ttfamily \textbf{PK} id : INT (AI)\\
            \scriptsize\ttfamily \textbf{UQ} codigo : VARCHAR(50)\\
            \scriptsize\ttfamily \phantom{XX} nombre : VARCHAR(150)\\
            \scriptsize\ttfamily \phantom{XX} sigla : VARCHAR(30)\\
            \scriptsize\ttfamily \textbf{FK} padre\_id : INT (NULL)\\
            \scriptsize\ttfamily \phantom{XX} nivel : INT\\
            \scriptsize\ttfamily \phantom{XX} activo : BOOLEAN
        \end{tabular}
    };

    % CONEXIONES Y CARDINALIDADES
    \draw[rel] (personas.south) -- (usuarios.north)
        node[card, near start, xshift=10pt] {1}
        node[card, near end, xshift=10pt] {0..1};

    \draw[rel] (usuarios.east) -- (userroles.west)
        node[card, near start, yshift=7pt] {1}
        node[card, near end, yshift=7pt] {N};

    \draw[rel] (roles.south) -- (userroles.north)
        node[card, near start, xshift=10pt] {1}
        node[card, near end, xshift=10pt] {N};

    \draw[rel] (ubicaciones.west) -- (userroles.east)
        node[card, near start, yshift=7pt] {1}
        node[card, near end, yshift=7pt] {N};

    \draw[rel] (ubicaciones.north) to[out=60,in=120,distance=1.8cm] 
        node[card, above] {Recursiva (1 : N) padre\_id} (ubicaciones.north);

\end{tikzpicture}
}
\caption{Diagrama Entidad-Relación (ERD) del Módulo de Identidad, Seguridad y Organigrama.}
\label{fig:erd_sprint1}
\end{figure}

% ==============================================================================
% SECCIÓN 3: SISTEMA DE SEGURIDAD Y AUTENTICACIÓN
% ==============================================================================
\section{Sistema de Seguridad, Autenticación y Autorización}

La seguridad del Sistema Wayka se diseñó bajo las pautas del estándar OWASP ASVS (\emph{Application Security Verification Standard}), implementando un esquema desacoplado de autenticación basada en tokens no intrusivos en estado (\emph{Stateless JWT}).

\subsection{Ciclo de Vida y Flujo Criptográfico del Token JWT}
El ciclo de autenticación opera de acuerdo con las siguientes fases:
\begin{enumerate}[leftmargin=1.5em]
    \item \textbf{Recepción de Credenciales:} El cliente transmite \texttt{login} y \texttt{password} cifrados vía canal seguro HTTPS (TLS 1.3).
    \item \textbf{Validación de Identidad y Hash:} El sistema consulta al usuario activo y realiza la comprobación computacional mediante la función unívoca \texttt{bcrypt.compare(password, password\_hash)}. Si la comparación es infructuosa, se responde de inmediato con código \texttt{401 Unauthorized}, sin divulgar si el error residió en el usuario o en la clave.
    \item \textbf{Carga del Contexto de Rol Activo:} Se evalúan los roles vigentes en la tabla \texttt{usuario\_roles}. Si el usuario cuenta con varios cargos, se activa automáticamente el rol designado como \texttt{es\_principal = 1}, permitiendo conmutar a otros roles mediante el endpoint \texttt{/api/auth/switch-role}.
    \item \textbf{Firma del Access Token y Refresh Token:} Se emiten dos artefactos criptográficos:
    \begin{itemize}
        \item \textbf{Access Token (Vigencia corta: 8 horas):} Contiene los \emph{claims} de sesión requeridos por los middlewares de autorización sin requerir lecturas continuas a disco:
        \[
        \text{Claims} = \{ \text{userId}, \text{username}, \text{roleCode}, \text{ubicacionOrgId}, \text{nivelAcceso}, \text{iat}, \text{exp} \}
        \]
        \item \textbf{Refresh Token (Vigencia extendida: 7 días):} Permite al cliente solicitar renovación transparente de credenciales.
    \end{itemize}
    \item \textbf{Intercepción por Middleware:} Cada solicitud a rutas protegidas es interceptada por \texttt{authMiddleware.js}, extrayendo el encabezado \texttt{Authorization: Bearer <token>} y adjuntando la información decodificada al objeto \texttt{req.user}.
\end{enumerate}

\subsection{Matriz de Permisos Basada en Roles (RBAC)}
El control de acceso gobierna las capacidades del sistema a través de la matriz de privilegios detallada en la Tabla~\ref{tab:rbac_matrix}.

\begin{table}[H]
\centering
\small
\caption{Matriz de Control de Acceso Basada en Roles (RBAC) -- Sprint 1}
\label{tab:rbac_matrix}
\begin{tabularx}{\linewidth}{@{}>{\raggedright\arraybackslash}p{4.7cm}*{4}{C}@{}}
\toprule
\textbf{Módulo / Funcionalidad del Sistema} & \textbf{Admin Sistema} & \textbf{Admin Wayka} & \textbf{Ventanilla Única} & \textbf{Funcionario} \\
\midrule
Gestión de Usuarios y Credenciales & \textbf{\color{colorExito}CRUD Total} & \color{colorAlerta}$\times$ Bloqueado & \color{colorAlerta}$\times$ Bloqueado & \color{colorAlerta}$\times$ Bloqueado \\
Gestión de Personas (Padrón Base) & \textbf{\color{colorExito}CRUD Total} & \color{colorAlerta}$\times$ Bloqueado & \color{colorAlerta}$\times$ Bloqueado & \color{colorAlerta}$\times$ Bloqueado \\
Gestión de Organigrama y Unidades & \textbf{\color{colorExito}CRUD Total} & \color{colorAlerta}$\times$ Bloqueado & \color{colorAlerta}$\times$ Bloqueado & \color{colorAlerta}$\times$ Bloqueado \\
Configuración de Parámetros Globales & \textbf{\color{colorExito}CRUD Total} & \color{colorAlerta}$\times$ Bloqueado & \color{colorAlerta}$\times$ Bloqueado & \color{colorAlerta}$\times$ Bloqueado \\
Administración de Tipos de Trámite & \color{colorAlerta}$\times$ Bloqueado & \textbf{\color{colorExito}Gestión Operativa} & \color{colorAlerta}$\times$ Bloqueado & \color{colorAlerta}$\times$ Bloqueado \\
Anulación / Habilitación de Trámites & \color{colorAlerta}$\times$ Bloqueado & \textbf{\color{colorExito}Exclusivo} & \color{colorAlerta}$\times$ Bloqueado & \color{colorAlerta}$\times$ Bloqueado \\
Redirección de Procesos en Conflicto & \color{colorAlerta}$\times$ Bloqueado & \textbf{\color{colorExito}Exclusivo} & \color{colorAlerta}$\times$ Bloqueado & \color{colorAlerta}$\times$ Bloqueado \\
Registro e Inicio de Trámites/VU & \color{colorAlerta}$\times$ Bloqueado & \color{colorAlerta}$\times$ Bloqueado & \textbf{\color{colorExito}Creación / Emisión} & \color{colorAlerta}$\times$ Bloqueado \\
Bloqueo / Desbloqueo Preventivo & \color{colorAlerta}$\times$ Bloqueado & \color{colorAlerta}$\times$ Bloqueado & \textbf{\color{colorExito}Operativo} & \color{colorAlerta}$\times$ Bloqueado \\
Bandeja Virtual (Escritorio de Trabajo) & \color{colorAlerta}$\times$ Bloqueado & \color{colorInfo}Supervisión & \textbf{\color{colorExito}Acceso Completo} & \textbf{\color{colorExito}Acceso Completo} \\
Atención de Trámites y Proveídos & \color{colorAlerta}$\times$ Bloqueado & \color{colorAlerta}$\times$ Bloqueado & \textbf{\color{colorExito}Habilitado} & \textbf{\color{colorExito}Habilitado} \\
Derivación Libre y Retroceso & \color{colorAlerta}$\times$ Bloqueado & \color{colorAlerta}$\times$ Bloqueado & \textbf{\color{colorExito}Habilitado} & \textbf{\color{colorExito}Habilitado} \\
Carga de Adjuntos Digitales (PDF) & \color{colorAlerta}$\times$ Bloqueado & \color{colorAlerta}$\times$ Bloqueado & \textbf{\color{colorExito}Habilitado} & \textbf{\color{colorExito}Habilitado} \\
Consulta Pública de Hojas de Ruta & \textbf{\color{colorExito}Permitido} & \textbf{\color{colorExito}Permitido} & \textbf{\color{colorExito}Permitido} & \textbf{\color{colorExito}Permitido} \\
Cambio de Contraseña (PIN Seguro) & \textbf{\color{colorExito}Propia} & \textbf{\color{colorExito}Propia} & \textbf{\color{colorExito}Propia} & \textbf{\color{colorExito}Propia} \\
\bottomrule
\end{tabularx}
\end{table}

% ==============================================================================
% SECCIÓN 4: ESPECIFICACIÓN DE LA API (ENDPOINTS DESARROLLADOS)
% ==============================================================================
\section{Especificación de la API (Endpoints Desarrollados)}

La capa de API REST se implementó bajo el framework Express en Node.js, estructurando las respuestas en formato JSON uniforme compuesto por los campos \texttt{\{success, message, data\}}.

\subsection{Catálogo de Rutas HTTP del Sprint 1}
En la Tabla~\ref{tab:endpoints} se presenta el compendio de rutas disponibles en la versión \texttt{1.0.0-sprint1}.

\begin{table}[H]
\centering
\small
\caption{Endpoints Desarrollados y Expuestos en el Sprint 1}
\label{tab:endpoints}
\begin{tabularx}{\linewidth}{@{}l>{\raggedright\arraybackslash}p{5.5cm}Y>{\raggedright\arraybackslash}p{2.7cm}@{}}
\toprule
\textbf{Método} & \textbf{Ruta de Acceso (Path)} & \textbf{Descripción Técnica} & \textbf{Nivel / Rol Mínimo} \\
\midrule
\texttt{GET} & \texttt{/api/} & Metadatos del microservicio y estado general. & Público \\
\texttt{GET} & \texttt{/api/health} & Verificación de salud y conectividad con MySQL. & Público \\
\texttt{POST} & \texttt{/api/auth/login} & Autenticación de usuarios y generación de JWT. & Público \\
\texttt{GET} & \texttt{/api/auth/profile} & Recupera datos de usuario y roles asignados. & Token Requerido \\
\texttt{POST} & \texttt{/api/auth/switch-role} & Conmuta el rol activo en sesión multi-rol. & Token Requerido \\
\texttt{POST} & \texttt{/api/auth/\allowbreak change-password} & Actualización de PIN con validación de complejidad. & Token Requerido \\
\bottomrule
\end{tabularx}
\end{table}

\subsection{Ejemplos de Cargas Útiles (Payloads JSON)}

\subsubsection{1. Inicio de Sesión Exitoso (\texttt{POST /api/auth/login})}
\begin{lstlisting}[language=json]
// Solicitud (Request)
{
  "login": "admin",
  "password": "admin123"
}

// Respuesta Exitosa (Response 200 OK)
{
  "success": true,
  "message": "Inicio de sesión exitoso",
  "data": {
    "token": "eyJhbGciOiJIUzI1NiIsInR5cCI6IkpXVCJ9...",
    "refreshToken": "eyJhbGciOiJIUzI1NiIsInR5cCI6IkpXVCJ9...",
    "user": {
      "id": 1,
      "username": "admin",
      "nombres": "Administrador",
      "apellidos": "Sistema Wayka",
      "ci": "1000001",
      "email": "admin@wayka.gob.bo",
      "cargo": "Administrador Principal",
      "activeRole": {
        "usuario_rol_id": 1,
        "rol_id": 1,
        "rol_codigo": "ADMIN_SISTEMA",
        "rol_nombre": "Administrador de Sistema",
        "ubicacion_org_id": 6,
        "ubicacion_codigo": "UNI-SIS",
        "ubicacion_nombre": "Unidad de Tecnologías de Información y Sistemas",
        "nivel_acceso": "CONTROL_TOTAL",
        "es_principal": 1
      }
    }
  }
}
\end{lstlisting}

\subsubsection{2. Conmutación Dinámica de Rol (\texttt{POST /api/auth/switch-role})}
\begin{lstlisting}[language=json]
// Solicitud (Request)
{
  "rolId": 2,
  "ubicacionOrgId": 6
}

// Respuesta (Response 200 OK - Nuevo Token Re-firmado)
{
  "success": true,
  "message": "Rol activo cambiado a Administrador de Wayka",
  "data": {
    "token": "eyJhbGciOiJIUzI1NiIsInR5cCI6IkpXVCJ9.reSignedToken...",
    "activeRole": {
      "rol_id": 2,
      "rol_codigo": "ADMIN_WAYKA",
      "rol_nombre": "Administrador de Wayka",
      "ubicacion_org_id": 6,
      "ubicacion_nombre": "Unidad de Tecnologías de Información y Sistemas"
    }
  }
}
\end{lstlisting}

% ==============================================================================
% SECCIÓN 5: PRUEBAS Y VALIDACIÓN (QUALITY ASSURANCE)
% ==============================================================================
\section{Pruebas y Validación (Quality Assurance)}

La garantía de calidad del Sprint 1 se abordó mediante una batería de pruebas automatizadas que se ejecutan sin intervención manual en el entorno del servidor Node.js y MySQL.

\subsection{Resultados del Banco de Pruebas de Autenticación (\texttt{testAuth.js})}
Se implementó un arnés de pruebas HTTP dinámico sobre Express para validar la robustez del subsistema criptográfico:
\begin{itemize}[leftmargin=1.5em]
    \item \textbf{Test 1 (Login Válido):} Generación de token JWT, resolución de claims y asignación de rol principal por defecto (\texttt{ADMIN\_SISTEMA}). $\rightarrow$ \textbf{\color{colorExito}Aprobado (200 OK)}.
    \item \textbf{Test 2 (Rechazo de Credenciales Inválidas):} Verificación de contraseñas erróneas y cuentas inexistentes. $\rightarrow$ \textbf{\color{colorExito}Aprobado (401 Unauthorized)}.
    \item {\raggedright\textbf{Test 3 (Protección de Rutas y Manipulación):} Acceso sin cabecera \texttt{Authorization} y con firma alterada. $\rightarrow$ \textbf{\color{colorExito}Aprobado (401 Denegado)}.\par}
    \item \textbf{Test 4 (Conmutación Multi-Rol):} Verificación de actualización del token tras conmutar a rol \texttt{ADMIN\_WAYKA}. $\rightarrow$ \textbf{\color{colorExito}Aprobado (200 OK)}.
    \item \textbf{Test 5 (Políticas de PIN Seguro):} Validación de reglas de complejidad (mínimo 6 caracteres, longitud máxima, combinación obligatoria de letras y números) y restauración de clave de prueba. $\rightarrow$ \textbf{\color{colorExito}Aprobado (400 Bad Request / 200 OK)}.
    \item \textbf{Test 6 (Cuentas Operativas Demo):} Verificación de inicio de sesión de Ventanilla Única (\texttt{mfernandez}) y Funcionario (\texttt{cmamani}). $\rightarrow$ \textbf{\color{colorExito}Aprobado (200 OK)}.
\end{itemize}

\subsection{Resultados de Integridad del Modelo de Datos (\texttt{testDataModel.js})}
La suite de persistencia evaluó los 11 modelos DAO con validación estricta de restricciones de negocio:
\begin{table}[H]
\centering
\small
\caption{Resultados de la Validación de Integridad de Persistencia}
\label{tab:test_datamodel}
\begin{tabularx}{\linewidth}{@{}>{\raggedright\arraybackslash}p{2.8cm}Y>{\raggedright\arraybackslash}p{3cm}>{\centering\arraybackslash}p{2.5cm}@{}}
\toprule
\textbf{Módulo Evaluado} & \textbf{Regla de Negocio / Restricción Verificada} & \textbf{Resultado} & \textbf{Estado} \\
\midrule
\texttt{Persona} & Registro con CI unívoco y búsqueda por padrón civil. & Insert/Select OK & \textbf{\color{colorExito}100\% Exitoso} \\
\texttt{UbicacionOrg} & Cálculo recursivo de nivel jerárquico ($N_{\text{padre}} + 1$). & Nivel 2 Automático & \textbf{\color{colorExito}100\% Exitoso} \\
\texttt{Rol} & Protección contra borrado con asignaciones vigentes. & Excepción capturada & \textbf{\color{colorExito}100\% Exitoso} \\
\texttt{Usuario} & Generación de salt y hash bcrypt en creación. & Hash verificado & \textbf{\color{colorExito}100\% Exitoso} \\
\texttt{UsuarioRol} & Desmarque automático de otros roles principales. & Exclusividad OK & \textbf{\color{colorExito}100\% Exitoso} \\
\texttt{Correlativo} & Generación secuencial atómica (\texttt{TP-1/2026, TP-2/2026}). & Cero colisiones & \textbf{\color{colorExito}100\% Exitoso} \\
\texttt{Tramite} & Creación en estado \texttt{EN\_ATENCION} y movimiento inicial. & Orden 1 (INICIO) & \textbf{\color{colorExito}100\% Exitoso} \\
\texttt{Movimiento} & \textbf{Inmutabilidad estricta:} Rechazo de \texttt{UPDATE} y \texttt{DELETE}. & Mutación rechazada & \textbf{\color{colorExito}100\% Exitoso} \\
\texttt{SoftDelete} & Bloqueo de eliminación de Persona con Usuario activo. & Integridad protegida & \textbf{\color{colorExito}100\% Exitoso} \\
\bottomrule
\end{tabularx}
\end{table}

% ==============================================================================
% SECCIÓN 6: ESTADO DE AVANCE Y PRÓXIMOS PASOS
% ==============================================================================
\section{Estado de Avance y Próximos Pasos}

\subsection{Tabla de Entregables Planificados vs. Ejecutados}
El Sprint 1 ha cerrado formalmente con un \textbf{100\% de cumplimiento} de acuerdo al cronograma establecido en el documento de requerimientos maestros.

\begin{table}[H]
\centering
\small
\caption{Balance de Cumplimiento del Sprint 1 (01 Sep -- 14 Sep 2026)}
\label{tab:balance_sprint1}
\begin{tabularx}{\linewidth}{@{}Ycc>{\raggedright\arraybackslash}p{3.5cm}@{}}
\toprule
\textbf{Objetivo Técnico / Tarea del Sprint} & \textbf{Planificado} & \textbf{Completado} & \textbf{Observaciones / Estado} \\
\midrule
Configuración de entorno Node.js, Express y MySQL 8.0+ & 100\% & 100\% & Entorno de desarrollo Docker/Local operativo. \\
Esquema Relacional unificado y 11 scripts de migración & 100\% & 100\% & Migraciones versionadas mediante \texttt{migrator.js}. \\
Módulo Criptográfico y Autenticación JWT Stateless & 100\% & 100\% & Sockets cerrados, Bearer token y claims de rol. \\
Matriz RBAC y Middleware de Autorización por Rol & 100\% & 100\% & Protección granular en API. \\
Carga de Semillas Base (Roles, Organigrama, Cuentas Demo) & 100\% & 100\% & 3 usuarios base y 6 unidades del organigrama. \\
Modelos DAO para Entidades Base (CRUD y Soft-delete) & 100\% & 100\% & Lógica de negocio e integridad en JavaScript. \\
\bottomrule
\end{tabularx}
\end{table}

\end{document}
